\chapter{ЭКСПЕРИМЕНТАЛЬНОЕ ИССЛЕДОВАНИЕ И АНАЛИЗ РЕЗУЛЬТАТОВ}

\section{Методика вычислительного эксперимента}

Чтобы оценить качество обучения более надёжно, была проведена серия экспериментов с разным числом эпизодов. Для каждого значения числа эпизодов агент обучался 5 раз с независимой случайной инициализацией. После этого каждая полученная Q-таблица оценивалась на 100 тестовых играх без исследования, то есть при выборе действий только по уже выученным оценкам. В качестве основного показателя использовался средний score на тестовых играх, а в качестве дополнительного показателя --- средний лучший score, достигнутый в серии оценивания.

\begin{table}[H]
    \centering
    \caption{Схема вычислительного эксперимента}
    \begin{tabularx}{\textwidth}{>{\raggedright\arraybackslash}p{0.34\textwidth} X}
        \toprule
        Характеристика & Описание \\
        \midrule
        Число независимых обучений & 5 запусков для каждого набора эпизодов \\
        Число тестовых игр & 100 игр для каждой обученной Q-таблицы \\
        Число эпизодов обучения & 50, 100, 200, 300, 500, 800, 1000, 1200 \\
        Основной показатель & Средний score на тестовых играх \\
        Дополнительный показатель & Средний лучший score в серии тестов \\
        Назначение эксперимента & Оценка зависимости качества игры от объёма обучающего опыта \\
        \bottomrule
    \end{tabularx}
\end{table}

\section{Результаты обучения при различном числе эпизодов}

\begin{table}[H]
    \centering
    \caption{Результаты обучения агента}
    \begin{tabular}{ccc}
        \toprule
        Эпизоды & Средний score & Средний лучший score \\
        \midrule
        50 & 4.262 & 15.800 \\
        100 & 7.544 & 19.200 \\
        200 & 9.760 & 19.400 \\
        300 & 9.912 & 18.800 \\
        500 & 13.958 & 21.400 \\
        800 & 15.422 & 21.200 \\
        1000 & 15.494 & 21.400 \\
        1200 & 16.230 & 21.600 \\
        \bottomrule
    \end{tabular}
\end{table}

\begin{figure}[H]
    \centering
    \begin{tikzpicture}
        \begin{axis}[
            width=0.92\textwidth,
            height=0.48\textwidth,
            xlabel={Число эпизодов обучения},
            ylabel={Значение метрики},
            xmin=50, xmax=1200,
            ymin=0, ymax=24,
            xtick={50,100,200,300,500,800,1000,1200},
            ymajorgrids=true,
            xmajorgrids=true,
            grid style={dashed,gray!30},
            legend style={at={(0.5,-0.2)}, anchor=north, legend columns=2, draw=none},
            tick label style={font=\small},
            label style={font=\small}
        ]
            \addplot[
                color=blue,
                mark=*,
                line width=1.1pt
            ] coordinates {
                (50,4.262)
                (100,7.544)
                (200,9.760)
                (300,9.912)
                (500,13.958)
                (800,15.422)
                (1000,15.494)
                (1200,16.230)
            };
            \addlegendentry{Средний score}

            \addplot[
                color=red!75!black,
                mark=square*,
                line width=1.1pt
            ] coordinates {
                (50,15.800)
                (100,19.200)
                (200,19.400)
                (300,18.800)
                (500,21.400)
                (800,21.200)
                (1000,21.400)
                (1200,21.600)
            };
            \addlegendentry{Средний лучший score}
        \end{axis}
    \end{tikzpicture}
    \caption{Зависимость среднего и лучшего результата от числа эпизодов обучения}
\end{figure}

График и таблица показывают устойчивый рост качества игры по мере накопления опыта. Наиболее заметное улучшение наблюдается в диапазоне от 50 до 500 эпизодов: средний score быстро увеличивается, а лучший результат выходит на уровень около 20 очков. После 800 эпизодов рост замедляется, что указывает на постепенный выход обучения на плато. Это означает, что агент действительно усваивает полезную стратегию, но дальнейшее увеличение числа эпизодов даёт уже не столь большой прирост качества.

\section{Интерпретация результатов обучения}

Для машинного обучения важен не только рост метрик, но и то, почему он происходит. В этом проекте агент постепенно заполняет Q-таблицу полезными оценками действий в часто встречающихся состояниях. Из-за этого на поздних этапах он реже выбирает опасные ходы и чаще движется так, чтобы сохранить безопасность или добраться до пищи.

Нужно учитывать, что обучение остаётся случайным процессом. Разные запуски могут давать немного разные результаты из-за случайного положения пищи и исследовательских действий агента. Но общий рост показателей всё равно хорошо заметен.

\section{Анализ влияния гиперпараметров на обучение}

Результат работы агента зависит не только от самого алгоритма, но и от выбора гиперпараметров. В данном проекте особенно важны коэффициент обучения $\alpha$, коэффициент дисконтирования $\gamma$, параметр исследования $\varepsilon$ и скорость его уменьшения. Даже при одной и той же игровой среде изменение этих величин может заметно повлиять на скорость обучения и на устойчивость итоговой стратегии.

Коэффициент $\alpha$ определяет, насколько сильно новые наблюдения влияют на уже накопленные оценки. Если сделать $\alpha$ слишком маленьким, агент будет учиться очень медленно: новые удачные и неудачные ситуации почти не будут менять таблицу. Если же $\alpha$ окажется слишком большим, оценки начнут резко колебаться, и обучение станет менее устойчивым. В текущей работе используется значение 0.1, которое даёт достаточно плавное, но заметное обновление знаний.

Коэффициент $\gamma$ показывает, насколько агент учитывает будущую награду. При слишком маленьком $\gamma$ агент будет ориентироваться почти только на ближайший шаг и хуже понимать, что безопасный манёвр сейчас может привести к пище через несколько ходов. При слишком большом значении обучение может становиться более чувствительным к далёким и неопределённым последствиям. Значение 0.9 в этой работе позволяет учитывать долгосрочную цель, но при этом не делает поведение чрезмерно инерционным.

Параметр $\varepsilon$ отвечает за баланс между исследованием и использованием уже найденных удачных действий. На старте обучения агенту важно часто пробовать разные ходы, поэтому начальное значение $\varepsilon$ равно 1.0. Затем этот параметр постепенно уменьшается. Если уменьшать его слишком быстро, агент может слишком рано закрепиться на слабой стратегии. Если уменьшать слишком медленно, обучение будет дольше оставаться случайным. Поэтому множитель $\texttt{EPSILON\_DECAY}=0.995$ и нижняя граница $\texttt{EPSILON\_MIN}=0.05$ здесь выбраны как разумный компромисс между поиском новых действий и закреплением уже найденных.

\begin{table}[H]
    \centering
    \caption{Роль гиперпараметров обучения}
    \begin{tabularx}{\textwidth}{>{\raggedright\arraybackslash}p{0.16\textwidth} >{\raggedright\arraybackslash}p{0.18\textwidth} X X}
        \toprule
        Параметр & Роль & Слишком малое значение & Слишком большое значение \\
        \midrule
        $\alpha$ & Скорость обновления Q-оценок & Обучение идёт медленно, агент слабо реагирует на новый опыт & Оценки сильно колеблются, обучение становится менее устойчивым \\
        $\gamma$ & Учёт будущей награды & Агент ориентируется почти только на ближайший шаг & Поведение может стать слишком зависимым от далёких последствий \\
        $\varepsilon$ & Доля случайных действий & Агент мало исследует среду и может рано застрять в слабой стратегии & Агент слишком долго действует почти случайно \\
        $\texttt{EPSILON\_DECAY}$ & Скорость уменьшения исследования & Исследование заканчивается слишком быстро & Исследование сохраняется слишком долго \\
        \bottomrule
    \end{tabularx}
\end{table}

\section{Ограничения признакового описания состояния}

Сильной стороной проекта является компактное описание состояния, но именно оно же задаёт и границы применимости метода. Вектор из 11 бинарных признаков хорошо работает в учебной постановке, потому что хранит только самую важную информацию: опасность рядом, текущее направление движения и относительное положение пищи. Этого достаточно, чтобы агент научился избегать простых столкновений и выбирать более разумные ходы.

Однако такое описание не передаёт полную картину игрового поля. Агент не знает точную форму собственного тела на дальних клетках, не видит длинные обходные маршруты и не получает детальной информации о будущих ловушках, которые могут возникнуть через несколько ходов. Поэтому при усложнении среды или при увеличении длины змейки такая модель может начать терять качество.

С практической точки зрения это важное наблюдение для курсовой работы. Оно показывает, что успех агента определяется не только формулой Q-learning, но и тем, насколько удачно подобраны признаки состояния. Если признаки слишком бедные, агент просто не сможет отличить по-настоящему хорошие ситуации от плохих. Если же признаки сделать слишком подробными, пространство состояний вырастет настолько, что табличный метод станет неудобным.

Таким образом, текущее представление состояния является компромиссным. Оно достаточно простое для табличного метода и достаточно информативное для учебной игры, но уже на этом уровне показывает, где начинаются ограничения такого подхода.

\section{Переход к более сложным методам: почему интересен DQN}

Когда пространство состояний становится слишком большим, табличный Q-learning начинает терять эффективность. Для каждой новой комбинации признаков нужно хранить отдельные значения, а число таких комбинаций быстро растёт. В более сложных задачах, где агент должен учитывать полную картину поля, такой способ хранения уже плохо масштабируется.

Именно поэтому логичным следующим шагом выглядит переход к DQN и другим методам глубинного обучения с подкреплением. В этом случае Q-функция хранится не в таблице, а аппроксимируется нейронной сетью. Это позволяет работать с более богатыми наблюдениями и переносить знания между похожими состояниями, а не учить каждую ситуацию почти отдельно.

Для данной курсовой работы переход к DQN не был основной целью, поскольку сама задача строилась как понятная и интерпретируемая учебная реализация. Тем не менее сравнение с DQN важно с методической точки зрения. Оно показывает, что выбранный в проекте табличный подход хорошо подходит для простых и компактных постановок, но при дальнейшем усложнении задачи более естественным становится использование нейросетевых моделей.

\begin{table}[H]
    \centering
    \caption{Сравнение табличного Q-learning и DQN}
    \begin{tabularx}{\textwidth}{>{\raggedright\arraybackslash}p{0.24\textwidth} X X}
        \toprule
        Критерий & Табличный Q-learning & DQN \\
        \midrule
        Представление Q-функции & Явная таблица значений & Нейронная сеть \\
        Подходящие состояния & Компактные и дискретные & Более сложные и высокоразмерные \\
        Интерпретируемость & Высокая, проще объяснять шаги обучения & Ниже, так как поведение задаётся параметрами сети \\
        Требования к ресурсам & Небольшие & Выше по памяти, времени и настройке \\
        Уместность для текущей работы & Очень высокая & Полезен как направление дальнейшего развития \\
        \bottomrule
    \end{tabularx}
\end{table}

\section{Ограничения метода и направления развития}

Главное ограничение текущего решения связано с тем, что здесь используется табличный Q-learning. Он хорошо работает, пока состояние описывается компактным набором признаков. Но если перейти к более сложным наблюдениям, например к полному изображению поля, число состояний станет слишком большим.

В дальнейшем можно подробнее изучить влияние параметров $\alpha$, $\gamma$ и $\varepsilon$, попробовать другое описание состояния, увеличить размер поля и провести дополнительные серии экспериментов с разными схемами награды. Это позволило бы лучше понять, какие именно проектные решения сильнее всего влияют на поведение агента.

Кроме того, перспективным направлением остаётся переход к более сложным методам, например DQN. Такой шаг был бы особенно полезен в тех задачах, где табличный метод уже не справляется из-за большого числа состояний или из-за необходимости работать с более подробным описанием игрового поля.
